% !TEX root =  main.tex
\chapter{Query Generation Approach}
\label{chap:query_generation}
\pagestyle{plain}

This chapter presents the detailed approach for generating diverse, realistic SQL query workloads using large language models. Query generation is the foundational stage of the pipeline, as the quality and diversity of generated queries directly impacts the effectiveness of downstream model training. We describe the prompt engineering strategy, the multi-layer validation pipeline, the format conversion process, and the synthetic fallback generator.


\section{LLM-Based Query Generation}

Traditional approaches to generating training workloads for learned database components fall into two broad categories: manual collection from query logs and synthetic generation using template-based or random sampling methods. Query logs require access to production systems, raise privacy concerns, and may not cover the full spectrum of query patterns needed for robust model training. Template-based generators, on the other hand, produce queries that are syntactically valid but often lack the structural variety and realistic predicate distributions found in real workloads. Random sampling can generate uniformly distributed queries but misses the complex correlations and patterns that characterize real-world usage.

Large language models offer a compelling alternative because they have been trained on vast corpora of natural language and code, including millions of SQL queries. This training enables them to generate queries that exhibit realistic structural patterns, meaningful predicate combinations, and diverse join topologies without explicit programming of these patterns. Our framework uses the Ollama platform to host and serve LLMs locally, ensuring data privacy and eliminating dependency on external API services. The default model is Llama~3.2, selected for its strong code generation capabilities and reasonable computational requirements. The system interacts with Ollama through its REST API, enabling seamless switching between different model sizes and families without code changes.


\section{Prompt Design}

The quality of LLM-generated queries critically depends on the prompt design. We employ a carefully structured prompt template that provides the model with sufficient context to generate valid, diverse queries while constraining output to the supported query fragment. The prompt consists of five hierarchically organized components, as illustrated in Figure~\ref{fig:prompt_structure}: a role assignment that establishes the LLM as a SQL workload generator, the complete database schema context including DDL statements with table names, columns, data types, and foreign key relationships, optional column statistics providing min/max values for numeric columns to enable realistic predicate generation, a set of strict generation rules, and an output format specification.

\begin{figure}[htbp]
\centering
\begin{tikzpicture}[
    block/.style={draw, rectangle, rounded corners=3pt, minimum width=13cm, minimum height=1.3cm, text centered, font=\small},
    arrow/.style={->, thick, >=stealth},
    node distance=1.5cm
]
    \node [block, fill=blue!15] (role) {\textbf{1. Role Assignment}: ``You are a SQL query workload generator for a database.''};
    \node [block, fill=green!15, below=of role] (schema) {\textbf{2. Schema Context}: Complete DDL with table names, columns, types, and relationships};
    \node [block, fill=yellow!15, below=of schema] (stats) {\textbf{3. Column Statistics}: Min/max values and cardinalities for numeric columns};
    \node [block, fill=orange!15, below=of stats] (rules) {\textbf{4. Generation Rules}: Strict constraints on query structure and complexity};
    \node [block, fill=red!15, below=of rules] (output) {\textbf{5. Output Format}: JSON array specification with example queries};

    \draw [arrow] (role) -- (schema);
    \draw [arrow] (schema) -- (stats);
    \draw [arrow] (stats) -- (rules);
    \draw [arrow] (rules) -- (output);

    \node [left=0.3cm of role, font=\scriptsize, text=gray, rotate=90, anchor=south] {Context};
    \node [left=0.3cm of rules, font=\scriptsize, text=gray, rotate=90, anchor=south] {Constraints};

\end{tikzpicture}
\caption{Structure of the LLM prompt template, consisting of five hierarchically organized components that guide the model from context understanding to constrained output generation.}
\label{fig:prompt_structure}
\end{figure}

The generation rules are critical for ensuring that the produced queries are compatible with the MSCN encoding. The prompt enforces that all queries use the \texttt{SELECT COUNT(*)} format exclusively, involve one to three tables from the provided schema with proper aliases, include join conditions using defined foreign key relationships, and use only simple comparison operators (\texttt{=}, \texttt{<}, \texttt{>}) with numeric values connected by \texttt{AND}. Constructs such as \texttt{OR}, \texttt{IN}, \texttt{LIKE}, \texttt{BETWEEN}, subqueries, and aggregations are explicitly forbidden. These constraints are necessary because the MSCN encoding assumes a specific query structure where predicates are AND-connected simple comparisons over numeric columns; queries violating these constraints would either fail during encoding or produce incorrect feature representations.

The prompt instructs the LLM to return queries as a JSON array of objects, each containing a single \texttt{"sql"} field. This structured output format enables reliable parsing and reduces the need for complex text extraction. A fallback regex-based extraction is used when the LLM includes additional text around the JSON.


\section{Diversity and Batch Generation}

To maximize workload diversity across generation batches, we employ a rotating set of diversity hints. Each batch receives a different textual hint that steers the LLM toward generating queries with specific characteristics, ranging from single-table queries with various filters to complex three-table joins with multiple predicates, and from point lookups to queries with extreme selectivities. The hint index is computed as $(\lfloor \text{existing\_count} / \text{batch\_size} \rfloor) \bmod 5$, ensuring systematic rotation through all diversity profiles as generation progresses. This approach avoids the mode collapse problem where the LLM repeatedly generates structurally similar queries.

Generating all queries in a single LLM call would exceed token limits and reduce diversity. Instead, the system generates queries in configurable batches, with the total number of batches determined by $\lceil \text{total\_queries} / \text{batch\_size} \rceil$. Each batch receives a fresh prompt with the current diversity hint, ensuring variety across the full workload. LLM generation is inherently stochastic and may fail due to network errors, malformed JSON output, or hallucination. Each batch is therefore retried up to three times with increasing delays between attempts, achieving near-complete batch success rates in practice. Algorithm~\ref{alg:batch_generation} presents the batch generation procedure.

\begin{algorithm}[htbp]
\caption{LLM-Based Query Generation}
\label{alg:batch_generation}
\begin{algorithmic}[1]
\State Determine how many batches are needed to reach the target number of queries
\For{each batch}
    \State Choose a diversity hint based on how many queries have been generated so far
    \State Construct a prompt containing the database schema, column statistics, and the chosen hint
    \For{up to 3 attempts}
        \State Send the prompt to the language model and receive a response
        \State Try to extract valid SQL queries from the response
        \If{at least one valid query was found}
            \State Add the extracted queries to the collection and move to the next batch
        \EndIf
    \EndFor
\EndFor
\State \Return the collected queries, trimmed to the target count
\end{algorithmic}
\end{algorithm}


\section{Multi-Layer Validation Pipeline}

Raw LLM output must undergo rigorous validation before entering the training pipeline. Despite careful prompt engineering, LLMs can produce queries that are syntactically invalid, reference non-existent schema elements, or use unsupported SQL constructs. Our validation pipeline applies four sequential filtering stages, each progressively more expensive but more thorough, as illustrated in Figure~\ref{fig:validation_pipeline}.

\begin{figure}[htbp]
\centering
\begin{tikzpicture}[
    node distance=1.8cm,
    block/.style={draw, rectangle, rounded corners=3pt, minimum width=5.5cm, minimum height=1.1cm, align=center, font=\small},
    reject/.style={draw, rectangle, rounded corners=3pt, minimum width=2.5cm, minimum height=0.8cm, align=center, font=\small, fill=red!15},
    arrow/.style={->, thick, >=stealth},
    rarrow/.style={->, thick, >=stealth, red!70},
]
    \node [block, fill=gray!15] (raw) {Raw LLM Output (N queries)};
    \node [block, fill=blue!15, below=of raw] (struct) {Layer 1: Structural Validation\\{\scriptsize Regex-based keyword filtering}};
    \node [block, fill=green!15, below=of struct] (parse) {Layer 2: SQL Parsing\\{\scriptsize FROM/WHERE clause extraction}};
    \node [block, fill=yellow!15, below=of parse] (schema) {Layer 3: Schema Validation\\{\scriptsize Table/column existence checks}};
    \node [block, fill=orange!15, below=of schema] (format) {Layer 4: Format Conversion\\{\scriptsize MSCN-compatible structure}};
    \node [block, fill=green!30, below=of format] (valid) {Valid Queries ($\leq$ N)};

    \draw [arrow] (raw) -- (struct);
    \draw [arrow] (struct) -- (parse);
    \draw [arrow] (parse) -- (schema);
    \draw [arrow] (schema) -- (format);
    \draw [arrow] (format) -- (valid);

    \node [reject, right=2.5cm of struct] (r1) {Reject: Forbidden\\keywords found};
    \node [reject, right=2.5cm of parse] (r2) {Reject: Parse\\failure};
    \node [reject, right=2.5cm of schema] (r3) {Reject: Unknown\\table/column};
    \node [reject, right=2.5cm of format] (r4) {Reject: Conversion\\failure};

    \draw [rarrow] (struct.east) -- (r1.west);
    \draw [rarrow] (parse.east) -- (r2.west);
    \draw [rarrow] (schema.east) -- (r3.west);
    \draw [rarrow] (format.east) -- (r4.west);

\end{tikzpicture}
\caption{Multi-layer validation pipeline for LLM-generated queries. Each layer progressively filters invalid queries, with rejected queries logged for analysis.}
\label{fig:validation_pipeline}
\end{figure}

The first layer is a fast, regex-based filter that rejects queries containing forbidden SQL patterns such as \texttt{OR}, \texttt{IN}, \texttt{LIKE}, \texttt{BETWEEN}, subqueries, and aggregations. This layer catches the most common LLM errors at minimal computational cost. The second layer parses each surviving query to extract its FROM and WHERE clauses, handling both comma-separated join syntax and explicit \texttt{JOIN...ON} syntax and producing a normalized representation of tables, join conditions, and filter predicates.

The third layer validates the parsed query against the database schema, verifying that all referenced table names, column names, and aliases actually exist. Join conditions are checked to confirm that they reference columns participating in defined foreign key relationships. The fourth and final layer converts valid queries from their parsed representation into the MSCN-compatible structured format, decomposing each query into three sets: tables (with aliases), joins (as equality conditions), and predicates (as column-operator-value tuples). This conversion includes operator standardization, value normalization, case normalization, and alias resolution to ensure consistent encoding downstream.


\section{Synthetic Query Generator}

For testing and environments where an LLM is unavailable, the framework includes a deterministic synthetic query generator. This generator produces structurally valid queries using parameterized random sampling based on hardcoded table definitions from the target schema. It uses weighted random selection for table count, biasing toward two-table queries to match typical workload distributions, and all multi-table queries include the central fact table to reflect the star-schema topology of the benchmark database.

Table~\ref{tab:gen_comparison} compares the characteristics of queries produced by the two generation approaches.

\begin{table}[htbp]
\centering
\caption{Comparison of LLM-based and synthetic query generation approaches.}
\label{tab:gen_comparison}
\begin{tabular}{|l|c|c|}
\hline
\textbf{Property} & \textbf{LLM-Based} & \textbf{Synthetic} \\
\hline
Schema adaptability & Automatic & Requires manual updates \\
Predicate diversity & High (learned patterns) & Medium (uniform random) \\
Join topology variety & High & Limited to star schema \\
Value distribution & Realistic (stat-aware) & Uniform within range \\
Validation pass rate & $\sim$95\% & 100\% (by construction) \\
Reproducibility & Non-deterministic & Fully deterministic \\
\hline
\end{tabular}
\end{table}

The LLM-based generator produces higher-quality workloads with greater diversity, while the synthetic generator provides a deterministic, zero-dependency fallback. In practice, the LLM-based generator is recommended for final experiments, while the synthetic generator serves well for rapid development and testing.


\section{Summary}

This chapter presented the detailed approach for generating SQL query workloads using large language models. The key contributions include a structured prompt engineering strategy that achieves high validation pass rates, a four-layer validation pipeline that ensures only well-formed and schema-compatible queries enter the training pipeline, a diversity hint rotation mechanism that prevents mode collapse in generated workloads, and a deterministic synthetic generator for testing environments. The next chapter describes how these generated queries are labeled, encoded, and used to train the MSCN cardinality estimator through active learning.
